%%%%%%%%%%  scan idx 511  |  folio 498  %%%%%%%%%%
\SGAPrintPage{498}%
\sgaStableTarget{sga6.e09.expose}\section*{\underline{EXPOSÉ IX}\\[3pt] \underline{\underline{SOME COMPUTATIONS OF GROUPS $K$}}}
\begin{center}
\underline{by P. BERTHELOT}
\end{center}

This exposé corresponds to no oral exposé, and will not be used in the rest of the seminar. It contains a number of computations of groups $K_\bullet$ for various fibrations: vector bundles, projective bundles, torsors under a linear group, etc\ldots{} Throughout this exposé, the base scheme will be supposed noetherian, and the notation $K_\bullet(X)$ will denote the Grothendieck group of coherent sheaves on a scheme $X$. Recall that if $X$ is noetherian, the homomorphism from the $K$ of the category of coherent $\underline{O}_X$-Modules to the $K$ of the category of pseudo-coherent complexes on $X$ is an isomorphism (IV 2.4).

The results presented here are analogous to those obtained for the Chow ring, which can be found in Grothendieck's exposés in the Chevalley Seminar 1958. The question also arises of determining to what extent statements of the same type can be obtained for the graded group associated with $K_\bullet$ for the filtration defined by the dimension of the support (cf. X). It does not seem that the answer is known, even for Proposition 1.1 (homotopy exact sequence), which is one of the properties most used in the computations that follow.

\sgaStableTarget{sga6.e09.sec.1}\subsection*{1. Vector bundles}

\sgaStableTarget{sga6.e09.prop.1-1}\textbf{Proposition 1.1. (Homotopy exact sequence).} \uline{Let} $X$ \uline{be a noetherian scheme}, $X'$ \uline{a closed subscheme of} $X$, $U$ \uline{the open complement of} $X'$ \uline{in} $X$, $i$ \uline{the immersion of} $X'$ \uline{into} $X$, \uline{and} $j$ \uline{the immersion of} $U$ \uline{into} $X$. \uline{Then the sequence}:

%%%%%%%%%%  scan idx 512  |  folio 499  %%%%%%%%%%
\SGAPrintPage{499}%
\[
K_\bullet(X')\xrightarrow{i_*}K_\bullet(X)\xrightarrow{j^*}K_\bullet(U)\longrightarrow 0
\]
\uline{is
exact}.

If $\underline{F}'$ is a coherent $\underline{O}_{X'}$-Module, $i_*(\underline{F}')$ is an $\underline{O}_X$-Module whose support is contained in $X'$; consequently its restriction to $U$ is zero, and $j^*\circ i_*(\cl(\underline{F}'))=0$, whence, since the elements of the form $\cl(\underline{F}')$ generate $K_\bullet(X')$, $j^*\circ i_*=0$. This gives a homomorphism
\[
 \rho:K_\bullet(X)/\operatorname{Im}K_\bullet(X')\longrightarrow K_\bullet(U).
\]
We shall define an inverse homomorphism to $\rho$. Let $\underline{G}$ be a coherent $\underline{O}_U$-Module. Since $X$ is noetherian, there exists a coherent $\underline{O}_X$-Module $\underline{F}$ whose restriction to $U$ is isomorphic to $\underline{G}$. If $\underline{F}'$ is a second coherent $\underline{O}_X$-Module whose restriction to $U$ is isomorphic to $\underline{G}$, then we have in $K_\bullet(X)$:
\[
 \cl(\underline{F})-\cl(\underline{F}')\in \operatorname{Im}K_\bullet(X').
\]
To see this, regard $\underline{G}$ as a subsheaf of $j^*(\underline{F}\oplus \underline{F}')$ by means of the diagonal map; $\underline{G}$ then extends to a subsheaf $\underline{F}''$ of $\underline{F}\oplus \underline{F}'$, and we are reduced to showing that
\[
 \cl(\underline{F})-\cl(\underline{F}'')\in \operatorname{Im}K_\bullet(X'),
\]
and the analogous assertion for $\underline{F}'$. Let then $\underline{H}$ and $\underline{K}$ be the kernel and cokernel of the composite homomorphism:
\[
\underline{F}''\longrightarrow \underline{F}\oplus \underline{F}'\longrightarrow \underline{F} ;
\]
this homomorphism reduces to the identity of $\underline{G}$ over $U$, and therefore $\underline{H}$ and $\underline{K}$ have support contained in $X'$. If $\underline{I}$ is the ideal of $X'$ in $X$, it follows that there exist integers $m$ and $n$ such that $\underline{I}^m\cdot \underline{H}=0$, $\underline{I}^n\cdot \underline{K}=0$ (EGA I 9.3.4). One may therefore express $\underline{H}$ and $\underline{K}$ as successive extensions of Modules

%%%%%%%%%%  scan idx 513  |  folio 500  %%%%%%%%%%
\SGAPrintPage{500}%
annihilated by $\underline{I}$, hence of $\underline{O}_{X'}$-Modules, and we have:
\[
\cl(\underline{H})\in \operatorname{Im}K_\bullet(X'),\qquad
\cl(\underline{K})\in \operatorname{Im}K_\bullet(X'),
\]
giving the result. We thus define a function on the category of coherent $\underline{O}_U$-Modules with values in $K_\bullet(X)/\operatorname{Im}K_\bullet(X')$ by assigning to every coherent $\underline{O}_U$-Module the class of one of its extensions to $X$.

This function is additive. Indeed, let:
\[
0\longrightarrow \underline{G}'\longrightarrow \underline{G}\longrightarrow \underline{G}''\longrightarrow 0
\]
be an exact sequence of coherent $\underline{O}_U$-Modules. Let $\underline{F}$ be an extension of $\underline{G}$ to $X$; then there exists a sub-Module $\underline{F}'$ of $\underline{F}$ extending $\underline{G}'$, and $\underline{F}/\underline{F}'$ extends $\underline{G}''$. It therefore factors through $K_\bullet(U)$, defining the required homomorphism $K_\bullet(U)\to K_\bullet(X)/\operatorname{Im}K_\bullet(X')$. One checks immediately that it is inverse to $\rho$.

\sgaStableTarget{sga6.e09.prop.1-2}\textbf{Proposition 1.2.} \uline{Let} $X,S$ \uline{be two noetherian schemes}, \uline{and} $f$ \uline{a} flat \uline{morphism from} $X$ \uline{to} $S$. \uline{Suppose given a family of elements} $(\xi_i)_{i\in I}$ \uline{of} $K^{\bullet}(X)$, \uline{and consider} \uline{the homomorphism}:
\[
\varphi:(K_\bullet(S))^{(I)}\longrightarrow K_\bullet(X)
\]
\uline{defined by}:
\[
\varphi((x_i)_{i\in I})=\sum_{i\in I}\xi_i\cdot f^*(x_i),
\]
\uline{where} $f^*$ \uline{is the} ``\uline{inverse-image}'' \uline{homomorphism from} $K_\bullet(S)$ \uline{to} $K_\bullet(X)$, \uline{defined by} IV 2.12, \uline{and where} $K_\bullet(X)$ \uline{is regarded as a} $K^{\bullet}(X)$\uline{-module by} IV 2.10. \uline{Suppose that for every point} $s\in S$, \uline{the images of the} $\xi_i$ \uline{in} $K_\bullet(X_s)$ \uline{under} \uline{the homomorphism}:
\[
K^{\bullet}(X)\longrightarrow K^{\bullet}(X_s)\longrightarrow K_\bullet(X_s)
\]
\uline{generate} $K_\bullet(X_s)$ \uline{as an abelian group}. \uline{Then} $\varphi$ \uline{is surjective}.

%%%%%%%%%%  scan idx 514  |  folio 501  %%%%%%%%%%
\SGAPrintPage{501}%
Let $S'$ be a closed subscheme of $S$, and $U$ the open complement of $S'$ in $S$. Denote by $X_U$ the restriction of $X$ over $U$, and by $X_{S'}$ the restriction of $X$ over $S'$. Finally, still denote by $\xi_i$ the images of the $\xi_i$ in $K^{\bullet}(X_{S'})$ and $K^{\bullet}(X_U)$, and denote by:
\[
\varphi_{S'}:K_\bullet(S')^{(I)}\to K_\bullet(X_{S'}),\qquad
\varphi_U:K_\bullet(U)^{(I)}\to K_\bullet(X_U)
\]
the homomorphisms defined in the same way as $\varphi$.

Consider the diagram
\sgaStableTarget{sga6.e09.eq.1-2-1}\[
\begin{tikzcd}[column sep=tiny]
K_\bullet(S')^{(I)} \arrow[r,"i_*"] \arrow[d,"\varphi_{S'}"'] & K_\bullet(S)^{(I)} \arrow[r,"j^*"] \arrow[d,"\varphi"] & K_\bullet(U)^{(I)} \arrow[r] \arrow[d,"\varphi_U"] & 0\\
K_\bullet(X_{S'}) \arrow[r,"u_*"] & K_\bullet(X) \arrow[r,"v^*"] & K_\bullet(X_U) \arrow[r] & 0
\end{tikzcd}\tag{1.2.1}
\]
where $i$ and $j$ denote the immersions of $S'$ and $U$ into $S$, and $u$ and $v$ the immersions of $X_{S'}$ and $X_U$ into $X$. Let us show that if $\varphi_{S'}$ and $\varphi_U$ are surjective, the same then holds for $\varphi$.

The two horizontal rows of the diagram are exact by 1.1. One then sees that to show that $\varphi$ is surjective, it suffices to show that the two squares commute.

Let $(x'_i)_{i\in I}\in K_\bullet(S')^{(I)}$; we have:
\[
\varphi\circ i_*((x'_i))=\varphi((i_*(x'_i)))=\sum_i \xi_i\cdot f^*(i_*(x'_i));
\]
\[
u_*\circ\varphi_{S'}((x'_i))=u_*\left(\sum_i \xi_i\cdot f^*_{S'}(x'_i)\right)=\sum_i \xi_i\cdot u_*(f^*_{S'}(x'_i)),
\]
the last equality following from the projection formula IV 2.11.1.2. The equality $\varphi\circ i_*=u_*\circ\varphi_{S'}$ therefore follows from the equality $f^*\circ i_*=u_*\circ f^*_{S'}$, which holds by IV 3.1.1 since $i$ is
proper and $f$ flat.

%%%%%%%%%%  scan idx 515  |  folio 502  %%%%%%%%%%
\SGAPrintPage{502}%
Let $(x_i)_{i\in I}\in K_\bullet(S)^{(I)}$; we have:
\[
\varphi_U\circ j^*((x_i))=\sum_i \xi_i\cdot f_U^*\circ j^*(x_i);
\]
\[
v^*\circ\varphi((x_i))=v^*\left(\sum_i \xi_i\cdot f^*(x_i)\right)=\sum_i \xi_i\cdot v^*\circ f^*(x_i)=\sum_i \xi_i\cdot f_U^*\circ j^*(x_i)
\]
the middle equality following from IV 2.12.2.

This result allows us to prove 1.2 by noetherian induction on the set of closed subschemes of $S$. If we suppose that for every closed subscheme $S'$ of $S$ distinct from $S$, $\varphi_{S'}$ is surjective, then the preceding argument shows that we may reduce to proving 1.2 for an irreducible affine base scheme, after replacing $S$ by an irreducible affine open subset. Moreover, the diagram
\[
\begin{tikzcd}
K_\bullet(S_{\red})^{(I)} \arrow[r] \arrow[d,"\varphi_{S_{\red}}"'] & K_\bullet(S)^{(I)} \arrow[d,"\varphi"]\\
K_\bullet(X_{S_{\red}}) \arrow[r] & K_\bullet(X),
\end{tikzcd}
\]
where $K_\bullet(X_{S_{\red}})\to K_\bullet(X)$ is surjective by 1.1, shows that we may also suppose $S$ reduced.

We therefore work over an integral affine base scheme. Put $S=\Spec(A)$, and denote by $K$ the field of fractions of $A$; if $s$ is the generic point of $S$, then $\underline{O}_{S,s}=K$.

Let $y\in K_\bullet(X)$. By hypothesis, $K_\bullet(X_s)$ is generated by the images of the $\xi_i$; if $y_s$ is the image of $y$ in $K_\bullet(X_s)$, we can find integers $n_i$ such that:
\[
y_s=\sum_i\xi_i\cdot n_i.
\]
